\section{Conclusion}
\label{sec:conclusion}

The OG--CLEWS link translates a solved resource strategy into the prices, investments,
health effects, and policy inputs that shape households, firms, government, and
different generations, and converts the resulting economic conditions back into inputs
for the next energy-model solve. Each relationship carries one mechanism, enters the
economy on one account, and can be inspected on its own, so a result can always be
traced back to the decision that produced it.

In the Philippine application, a decarbonisation
programme built around a coal moratorium from 2027 reshapes the power system while
leaving the economy moderately smaller in the long run: steady-state output changes by
$\cplYss\%$, with the adjustment carried mainly by wages ($\cplWss\%$) rather than
employment, and consumption by $\cplCss\%$. The same run converts a $\cplPMPct\%$
change in particulate exposure into roughly \cplDeaths{} avoided deaths and about
\cplMorbFTE{} full-time-equivalent workers per year of reduced illness. One
qualification travels with these numbers: nearly all of the output effect arrives
through the assumption that higher system costs pass into industry production costs,
and under the alternative productivity-based routing of the same cost signal the
aggregate sign reverses. The distributional account does not follow the familiar
budget-share intuition --- energy is a similar share of spending across income groups
here, so gains and losses are decided by earnings, assets, and transfers.

Three steps separate this system from application-grade policy numbers: a calibrated
monetary bridge, so that carbon and investment flows carry trusted units; a repeated
two-way solution, for which the mechanical seam is already tested; and a
country-calibrated health translation. Extending the same architecture to land and
water requires new resource-specific relationships and evidence, not a new method.
